\section{A border-valid dimension boundary for the mixed construction}

\subsection{O-017: the mixed staircase reaches the universal cut-family exponent}
Fix integers $k\geq2$, $s\geq1$, and a real exponent
$c\geq1$. Put $r=m^2$, and suppose $n=n(m)\geq m$
with $n(m)=O(m^c)$. Take one padding variable $L$,
$k-2$ early variables, and $R=r-k+1$ later variables.
For sufficiently large $m$, consider the whole linear space
of degree-$m$ forms having total later-variable degree at
most $s$. Its dimension is
\[
 D_m(k,s)=\sum_{h=0}^s
       \binom{R+h-1}{h}\binom{m-h+k-2}{k-2}
        =\frac{m^{2s+k-2}}{s!(k-2)!}
                                +O_{k,s}(m^{2s+k-3}).    \tag{O-017.1}
\]
If
\[
                            k+2s-4>2c,                 \tag{O-017.2}
\]
then not every form in that space, after multiplication by
$L^{n-m}$, belongs to $\Omega_n$ for all sufficiently large
$m$. This is a dimension statement about the actual orbit
closure, and therefore allows border representations.

\begin{proof}
For later degree $h$, the first binomial coefficient in
(O-017.1) counts the later monomials. The remaining degree
$m-h$ is distributed among the $k-2$ early variables and
$L$, giving the second binomial coefficient. Distinct
monomials are linearly independent, so the sum is the exact
dimension. Here $m\geq s$ is harmless because the conclusion
is asymptotic with fixed $s$. Since $R=m^2-k+1$,
the term $h=s$ has the displayed positive leading coefficient
and degree $2s+k-2$ in $m$; all smaller $h$ have lower
degree. Multiplication by the fixed nonzero monomial
$L^{n-m}$ is injective, so it preserves that dimension.

O-010 gives the dimension bound
\[
 \dim\bigl(\Omega_n\cap\operatorname{Sym}^n((\mathbb C^r)^*)\bigr)
                              \leq (r-2)n^2+2.         \tag{O-017.3}
\]
Its use here does not replace the orbit closure by the exact
image of substitutions. Indeed the intersection is the
\emph{closure} of the $r$-variable determinant-pencil image:
the projection to the first $r$ variables sends the full
determinant substitution image to that pencil image, and
continuity sends its closure into the closure of the latter.
A point of the intersection is fixed by this projection.
The reverse inclusion follows because each pencil is an
endomorphism substitution of the determinant. O-010 bounds
the dimension of this image closure by its explicit generic
gauge parametrization.

If the whole padded linear family were contained in
$\Omega_n$, its dimension would be at most the right-hand
side of (O-017.3). Since $r=m^2$ and $n=O(m^c)$, that
right-hand side is $O(m^{2c+2})$. Condition (O-017.2)
is exactly $2s+k-2>2c+2$. The positive leading term in
(O-017.1) then contradicts the required dimension inequality
for all sufficiently large $m$.
\end{proof}

\paragraph{Meaning for the staircase.}
The BIP staircase in F-017 uses $c=25$ and
$k=54-2s$, so $k+2s-4=50=2c$. The IP staircase
uses $c=4$, $k=12-2s$, again giving equality; a fixed
factor such as $n=3m^4+1$ does not change the exponent.
Thus those constructions reach the allowed exponent for
their universal cut families. Increasing the permitted later
degree by one at the same cut crosses (O-017.2).
No construction proving containment of that \emph{whole}
larger truncated family in the determinant orbit closure can
work at the same power range for all large $m$.

For a finite illustration in the IP range, take $m=100$,
$r=10000$, $n=3m^4+1$, $k=2$, and the raised cap $s=6$.
Then
\[
 D_m(2,6)=\binom{10005}{6}>(9998)(300000001)^2+2.
\]
The exact comparison is in the accompanying certificate.
For a BIP illustration, take $m=10^{15}$, $n=m^{25}$,
$k=4$, and the raised cap $s=26$; the same certificate
checks $D_m(4,26)>(m^2-2)m^{50}+2$ by integer arithmetic.

This boundary does not apply to a smaller family selected
by additional nested row bounds, nor to the actual family
of projected permanents. The dimension of an unrestricted
coefficient family is not the dimension or evaluation rank
of those smaller families. In particular, (O-017.2) does
not exhibit an actual multiplicity obstruction, and it does
not say that the F-017 staircase is a complete classification.
It identifies the limitation of extending the same universal
cut-family containment argument by one unit in the suffix cap.

\paragraph{Geometric dictionary.}
The truncated family is a finite torus-weight affine linear
subscheme, with the grading interpretation of SGA~3, I,
Proposition~4.7.3, pp.~35--36. Padding by a fixed section
is a closed linear embedding. The determinant side is a
scheme-theoretic image closure after base change to
$\mathbb C$, and its dimension is bounded by the explicit
parametrization in O-010. This comparison already includes
limits and makes no normality, cohomological purity, or
saturation assumption. It is separate from O-016's stronger
fixed-base-degree exponent statement, whose method requires
exact determinant representations.
